\section{Computability and the finite form}\label{sec:eff}

This section proves Theorem~\ref{thm:main}(iii). The output is a rational polynomial multiplied by a fixed
smooth spatial cutoff, and it is found by a search over finite certificates of trajectory containment: a
candidate polynomial is accepted when Euler trajectories from a grid covering the source balls, with an error
bound valid on the whole time interval, stay in the region where the cutoff equals one and end inside the target. Every
accepted candidate is correct (\Sref{Theorem}{thm:certificate-sound}). Termination needs one smooth compactly
supported Hamiltonian satisfying the packing specification with a positive margin. We take
$H[n,k,C,c,s]$ and compile it by Bernstein approximation (\S\ref{sec:eff-finite}); any packing Hamiltonian with
a margin would serve equally.

For the Epoch family the termination witness is Corollary~\ref{thm:main-epoch}, at density $1-\eps/2$; the
finite verifier is the same in every case.

A name for a real $x$ is a rational stream $q(p)$ with $|q(p)-x|\le2^{-p}$; a rational is named by its constant
stream. The finite verifier requires names for the source radii $r_i=\sqrt{c_i/\pi}$ and
target radius $R=\sqrt{C/\pi}$; these are computed from the supplied capacity
names by the programs in \Ssec{sec:efd}.
The symbol $|\cdot|$ denotes Euclidean norm, and $S$ is the flat step of
\Sref{Definition}{def:pre-flat}.

\begin{table}[ht]
\centering\footnotesize
\begin{tabular}{@{}L{.18\textwidth}L{.25\textwidth}L{.33\textwidth}L{.16\textwidth}@{}}
\toprule
Object & Definition & Effective data & Reference\\
\midrule
Input radii & $r_i=\sqrt{c_i/\pi}$, $R=\sqrt{C/\pi}$ & names with error $2^{-p}$; lower bound $1<\pi$ & \Ssec{sec:efd}\\
Source enclosures & $s_i^p$, $r^+_{i,p}$ & whole closed source covered; retained nodes within $(2d+3)2^{-p}+d\rho$ of a source & \Sref{Eq.}{eq:certificate-enclosures}\\
Polynomial field & $J_0\nabla P_e$, arguments clamped to $[0,1]\times[-\frac B2,\frac B2]^d$ & exact rational evaluation; integer $\Theta=h_B(e)$ bounding speed and Lipschitz constants & \Sref{Lem.}{lem:certificate-field}\\
Covering grid & mesh $\rho=2Q2^{-m}$ in $[-Q,Q]^d$ & $(2^m+1)^d$ nodes; maximum-norm radius $\rho$ & \Sref{Def.}{def:certificate-checks}\\
Euler table & $x_{j+1}=x_j+F_e(j/N,x_j)/N$ & error $E=(\rho+(\Theta+\Theta^2)/N)4^\Theta+\Theta/N$ on $[0,1]$ & \Sref{Lem.}{lem:certificate-euler}\\
Finite output & $\beta_BP_e$ & margins $B/2-|x_{j,l}|-E$ to the region where $\beta_B=1$ and $R^-_p-|x_N|-dE$ to the target boundary & \Sref{Thm.}{thm:certificate-sound}\\
Bernstein witness & compilation datum (C1)--(C7) & derivative envelope $M_0$; field error $\eta$; Gr\"onwall factor $(e^{L_X}-1)/L_X$ & Thm.~\ref{thm:eff-finite}\\
\bottomrule
\end{tabular}
\caption{Data for the finite verifier and its nonemptiness proof; $d=2n$.}
\label{tab:eff-data}
\end{table}

\begin{table}[ht]
\centering\footnotesize
\begin{tabular}{@{}L{.26\textwidth}L{.34\textwidth}L{.31\textwidth}@{}}
\toprule
Search & Nonempty because & Checkable condition\\
\midrule
Square-root names & none: a direct total program (integer square root at doubled precision, \Sref{Lemma}{lem:efd-roots-direct}) & ---\\
Candidate $(B,e,Q,p,m,N)$ & smaller-target geometric packing, compilation and source-neighbourhood margins & finite strict rational tests\\
First passing candidate & \Sref{Thm.}{thm:certificate-exists} & zero natural failure count\\
\bottomrule
\end{tabular}
\caption{Searches of the finite algorithm. The model parameters and their margins are given in the geometric sections.}
\label{tab:eff-searches}
\end{table}

\subsection{The finite form}\label{sec:eff-finite}

One estimate on the vector fields controls the time-one maps at all starting points at once.

\begin{lemma}[Gr\"onwall]\label{lem:eff-gronwall}
\uses{S-lem:cpl-gronwall-abstract}
Let $G,H\colon\R\times\R^{2n}\to\R$ be $C^\infty$, let $\eta>0$ and $L_X>0$, and suppose that
\begin{enumerate}[label=(\roman*)]
\item $|X_G(t,z)-X_H(t,z)|\le\eta$ for all $t\in[0,1]$ and all $z\in\R^{2n}$;
\item $X_H(t,\cdot)$ is $L_X$-Lipschitz on all of $\R^{2n}$ for every $t\in[0,1]$.
\end{enumerate}
If $\gamma_G,\gamma_H\colon[0,1]\to\R^{2n}$ are integral curves of $X_G$ and $X_H$ with
$\gamma_G(0)=\gamma_H(0)$, then $|\gamma_G(1)-\gamma_H(1)|\le\eta(e^{L_X}-1)/L_X$. If moreover
$\eta=b_Re^{-L_X}/2$ with $b_R>0$, then $\eta(e^{L_X}-1)/L_X<b_R$.
\end{lemma}

\begin{proof}
The difference $u=\gamma_G-\gamma_H$ satisfies $u(0)=0$ and
$|\dot u|\le|X_G(\gamma_G)-X_H(\gamma_G)|+|X_H(\gamma_G)-X_H(\gamma_H)|\le\eta+L_X|u|$, and Gr\"onwall's
inequality gives the first bound. With $\eta=b_Re^{-L_X}/2$ it equals $\tfrac{b_R}2(1-e^{-L_X})/L_X<b_R/2$, since
$1-e^{-x}<x$ for $x>0$. (\Sref{Lemma}{lem:cpl-gronwall} gives the same argument with all constants.)
\end{proof}

Both hypotheses must hold on the whole space: before the estimate is proved, the approximate trajectory is not
known to stay near the exact one, so a bound valid only near the exact trajectories could not be applied to it.
For smooth compactly supported $H$ such global bounds are finite: since $X_H=J_0\nabla_zH$ with $J_0$ orthogonal,
(ii) holds for every $L_X>0$ with $L_X\ge\max_{[0,1]\times\R^{2n}}\|D^2_zH\|$.

We now fix the compilation. Let $n\ge2$, $k\ge1$, $C>0$, $c\in(0,\infty)^k$ and $s\in(\R^{2n})^k$, and let $H$
satisfy $\Pack(n,k,C,c,s)$ (Definition~\ref{def:intro-spec}); thus $H$ is $C^\infty$ and vanishes outside
$\R\times K_H$ for a compact set $K_H\subset\R^{2n}$, its flow $\phi^t_H$ is global, and
$\phi^1_H(B(s_i,c_i))\subset\Bo(C)$ for every $i$. Write $\zeta_1,\dots,\zeta_{2n}$ for the real coordinates of
$z\in\R^{2n}$. A \emph{compilation datum} for $H$ consists of the following numbers.
\begin{enumerate}[label=(C\arabic*)]
\item A capacity $C_0\in(0,C)$ with $\phi^1_H(B(s_i,c_i))\subset B(C_0)$ for all $i$. Put
 $R_*=\sqrt{C_0/\pi}<R_C=\sqrt{C/\pi}$.
\item A rational $a_{\rm box}>0$ with $H(t,z)=0$ whenever $\max_l|\zeta_l|\ge a_{\rm box}$, and an integer
 $B>\max(2n,2a_{\rm box})$. The box cutoff is
 \[
 \beta_B(z)=\prod_{l=1}^{2n}\beta_{B,1}(\zeta_l),\qquad
 \beta_{B,1}(x)=1-S\Bigl(\frac{x^2-(B/2)^2}{(3B/4)^2-(B/2)^2}\Bigr),
 \]
 where $S$ is the flat step, so that $0\le\beta_B\le1$, $\beta_B=1$ on $[-B/2,B/2]^{2n}$, $\beta_B=0$ outside
 $(-3B/4,3B/4)^{2n}$, and $\beta_BH=H$.
\item Reals $L_X\ge1+\sup_{[0,1]\times\R^{2n}}\|D^2_zH\|$ and $C_\beta\ge1+\sup_{\R^{2n}}|\nabla\beta_B|$.
\item A rational $b_R$ with $0<b_R<(R_C-R_*)/4$; put $\eta=b_Re^{-L_X}/2$.
\item A real $M_0\ge1$ bounding all first and second partial derivatives on $[0,1]^{2n+1}$ of
 $\tilde H(\xi)=H(\xi_0,2B\xi_1-B,\dots,2B\xi_{2n}-B)$.
\item A rational $\delta$ with $0<\delta<\min\bigl(\tfrac12,\eta/(16(1+C_\beta))\bigr)$ and an integer
 $N_{\rm B}>\max\bigl(1,((2n+1)M_0/\delta)^2,4M_0/\delta\bigr)$.
\item For every $m\in\{0,\dots,N_{\rm B}\}^{2n+1}$ a rational $\tilde H_m$ with
 $|\tilde H_m-\tilde H(m/N_{\rm B})|<\delta/(2N_{\rm B}+1)$.
\end{enumerate}
With $b_{N,j}(x)=\binom Njx^j(1-x)^{N-j}$ put
\[
 P(t,z)=\sum_{m\in\{0,\dots,N_{\rm B}\}^{2n+1}}\tilde H_m\,b_{N_{\rm B},m_0}(t)\prod_{l=1}^{2n}b_{N_{\rm B},m_l}\Bigl(\frac{\zeta_l+B}{2B}\Bigr),
 \qquad \Hfin=\beta_B\cdot P.
\]
Since $B\in\N$ and $\tilde H_m\in\Q$, $P\in\Q[t,z]$, of degree at most $N_{\rm B}$ in each variable. Compilation
data exist: $C_0$ because $\bigcup_i\phi^1_H(B(s_i,c_i))$ is a compact subset of $\Bo(C)$, $a_{\rm box}$ because $H$
vanishes outside $\R\times K_H$, the bounds of (C3) and (C5) because they are maxima of continuous functions on
compact sets, and the rest because each is a strict lower or upper bound on a single number.

\begin{theorem}[Finite output]\label{thm:eff-finite}
\uses{lem:eff-gronwall, def:intro-spec, S-thm:cpl-final, S-prop:cpl-bernstein, S-lem:cpl-cutoff}
Let $H$ satisfy $\Pack(n,k,C,c,s)$. For every compilation datum (C1)--(C7) for $H$, the function
$\Hfin=\beta_B\cdot P$ satisfies $\Pack(n,k,C,c,s)$. More precisely: $\Hfin$ is $C^\infty$ and $\Hfin(t,z)=0$ for
all $t$ whenever $z\notin(-3B/4,3B/4)^{2n}$; $|X_{\Hfin}-X_H|<\eta$ on $[0,1]\times\R^{2n}$; and for every
$z\in\bigcup_iB(s_i,c_i)$,
\[
 |\phi^1_{\Hfin}(z)-\phi^1_H(z)|<b_R,\qquad |\phi^1_{\Hfin}(z)|<R_*+b_R<R_C,
\]
so that $\phi^1_{\Hfin}$ maps every $B(s_i,c_i)$ into $\Bo(C)$. If the source balls are pairwise disjoint,
their images are pairwise disjoint.
\end{theorem}

\begin{proof}[Proof (complete version: \Sref{Theorem}{thm:cpl-final})]
Write $\mathcal B_{N_{\rm B}}$ for the tensor Bernstein operator on $[0,1]^{2n+1}$. It satisfies
$\|\mathcal B_{N_{\rm B}}\tilde H-\tilde H\|_\infty\le(2n+1)M_0/(2\sqrt{N_{\rm B}})<\delta/2$ and
$\|\partial_j\mathcal B_{N_{\rm B}}\tilde H-\partial_j\tilde H\|_\infty\le
M_0\bigl((2n+1)/(2\sqrt{N_{\rm B}})+2/N_{\rm B}\bigr)<\delta$, by the binomial variance identity and summation by
parts (\Sref{Proposition}{prop:cpl-bernstein}). Because the Bernstein weights are nonnegative with sum $1$,
rounding changes values by less than $\delta/(2N_{\rm B}+1)$ and first derivatives by less than
$2N_{\rm B}\delta/(2N_{\rm B}+1)<\delta$. Hence on $[0,1]\times[-B,B]^{2n}$ we have $|P-H|<\delta$ and, since
$\partial_{\zeta_l}=(2B)^{-1}\partial_{\xi_l}$ and $B>2n$, $|\nabla_zP-\nabla_zH|\le\sqrt{2n}\,\delta/B<\delta$
(\Sref{Lemma}{lem:cpl-poly}). There, differentiating $\beta_BH=H$ gives
\[
 \nabla_z\Hfin-\nabla_zH=\beta_B\nabla_z(P-H)+(P-H)\nabla\beta_B ,
\]
whose norm is less than $\delta(1+C_\beta)<\eta/16$ by (C6); outside $(-B,B)^{2n}$ both $\Hfin$ and $H$ vanish on a
neighbourhood. As $X_G=J_0\nabla_zG$ with $J_0$ orthogonal, $|X_{\Hfin}-X_H|<\eta$ everywhere on
$[0,1]\times\R^{2n}$, and $X_H(t,\cdot)$ is $L_X$-Lipschitz by (C3). The field $X_{\Hfin}$ is smooth and
vanishes outside $\R\times[-3B/4,3B/4]^{2n}$, so its integral curves exist on $\R$ and are unique; the same
holds for $X_H$, which vanishes outside $\R\times K_H$. Lemma~\ref{lem:eff-gronwall} gives
$|\phi^1_{\Hfin}(z)-\phi^1_H(z)|<b_R$ for every $z$. For $z\in B(s_i,c_i)$ we have $|\phi^1_H(z)|\le R_*$ by
(C1), hence $|\phi^1_{\Hfin}(z)|<R_*+b_R<R_C$, that is $\pi|\phi^1_{\Hfin}(z)|^2<C$. Hence $\Hfin$ satisfies both
conditions of Definition~\ref{def:intro-spec}: it is $C^\infty$ and vanishes for $z\notin(-3B/4,3B/4)^{2n}$ and
all $t$, because $\beta_B$ does not depend on $t$; and every integral curve of $X_{\Hfin}$ on $[0,1]$ starting at
$z\in B(s_i,c_i)$ is the flow line, which ends in $\Bo(C)$. The images are disjoint because $\phi^1_{\Hfin}$ is
injective.
\end{proof}

The theorem applies to every compilation datum; the existence of one is used in the proof that the search
below terminates.

\subsection{The algorithm}\label{sec:eff-algorithm}

\begin{remark}[Accepted outputs]\label{rem:eff-certified}
The polynomial returned by the algorithm below need not be the Bernstein approximation of
$H[n,k,C,c,s]$ for any compilation datum. Bernstein compilation shows that some candidate is
accepted; every accepted candidate is correct by the soundness theorem.
\end{remark}

Put $d=2n$. A candidate consists of integers $B,Q,N\ge1$, precisions $p,m\ge0$,
and a finite rational arithmetic expression $e$ for a polynomial $P_e(t,z)$.
Let $\Theta=h_B(e)$ be the integer of \Sref{Definition}{def:certificate-program}.
It bounds the speed and the temporal and spatial Lipschitz constants, in maximum
norm, of the field with clamped arguments
\[
 F_e(t,z)=J_0\nabla_zP_e(\operatorname{clip}_{[0,1]}t,
                  \operatorname{clip}_{[-B/2,B/2]^d}z),
\]
where $\operatorname{clip}_X$ is the nearest-point projection to the box $X$.
Read the input names to error $2^{-p}$. From them form rational outer balls
$B(s_i^p,r^+_{i,p})$ containing the original closed sources, and a rational
$R^-_p\le\sqrt{C/\pi}$, as in \Sref{Equation}{eq:certificate-enclosures}.
First test $|s^p_{i,l}|+r^+_{i,p}<Q$ for every $i,l$.

Use the grid of mesh $\rho=2Q2^{-m}$ in $[-Q,Q]^d$. For each node $a$ and source
$i$, discard the node if $|a-s_i^p|^2>(r^+_{i,p}+d\rho)^2$.
Otherwise compute the exact rational Euler table
\[
 x_0=a,\quad x_{j+1}=x_j+N^{-1}F_e(j/N,x_j),\quad
 E=\left(\rho+\frac{\Theta+\Theta^2}{N}\right)4^\Theta+\frac{\Theta}{N}.
\]
Require $|x_{j,l}|+E<B/2$ for every $j\le N,l$, and
\[
 R^-_p-dE>0,\qquad |x_N|^2<(R^-_p-dE)^2.
\]
These are finite strict rational tests. Enumerate all candidate tuples and return
the first one passing every test; the output is $B$ and the program $e$, which determines $P_e\in\Q[t,z]$.

The real inputs are supplied as oracles and $n,k$ as ordinary input; \Ssec{sec:alg} gives the encoding.

\begin{theorem}[Termination and correctness]\label{thm:eff-algorithm}
There is one oracle machine with the following property. Let $n\ge2$, $k\ge1$, $C$, $c$ and $s$ satisfy the
hypotheses of Theorem~\ref{thm:main}, and supply as oracles names of $C,c_1,\dots,c_k$ and of the $2nk$ coordinates
of $s_1,\dots,s_k$. On input $(n,k)$ the machine halts and outputs an integer $B$ and a polynomial
$P\in\Q[t,z_1,\ldots,z_{2n}]$ such that $\beta_BP$ satisfies $\Pack(n,k,C,c,s)$. The output may depend on the names,
not only on the numbers they name. If the names are computable, for instance the constant names of rational
inputs, replacing the oracles by programs for them gives an ordinary algorithm.

There is also one oracle machine which, given a name of $\eps\in(0,1)$ as oracle and an integer $k\ge9$ as input,
computes names of the source radius $1$ and of the target radius $R=(k/(1-\eps/2))^{1/4}$ for the data
\[
 n=2,\quad C=\pi R^2,\quad c_i=\pi,\quad s_i=(3i-3,0,0,0),
\]
runs the same certificate search on them, and halts with $(B,P)$ such that $\beta_BP$ satisfies $\Pack(2,k,C,c,s)$;
here $k/R^4=1-\eps/2>1-\eps$.
\leanmark{original\_name\_algorithm, epoch\_name\_algorithm, epoch\_name\_algorithm\_packing}
\end{theorem}

\begin{proof}
For an accepted candidate, each source point lies within $\rho$ in maximum norm
of a grid node that cannot be discarded. The Euler error is at most $E$ throughout
$[0,1]$. The plateau tests therefore keep the trajectory of $F_e$ from that point in the open
box where $F_e=X_{\beta_BP_e}$, and the terminal test places its endpoint in the
open target. Smoothness, spatial compact support and uniqueness of integral curves give the
packing specification (\Sref{Theorem}{thm:certificate-sound}).

For termination, $H=H[n,k,C,c,s]$ maps the compact source balls into $\Bo(C)$, so their images lie in $\Bo(C_0)$
for some $C_0<C$; put $C_1=(C_0+C)/2$. Then $H$ satisfies $\Pack(n,k,C_1,c,s)$, and Theorem~\ref{thm:eff-finite} with
target $C_1$ and $B$ large gives a rational polynomial whose trajectories from the sources stay a positive distance
inside the region where $\beta_B=1$ and end in $\Bo(C_1)$, at a positive distance from the boundary of $\Bo(C)$.
These margins persist for starting points in a neighbourhood of the sources.
A node not discarded lies at maximum-norm distance at most
$(2d+3)2^{-p}+d\rho$ from an original source ball. Taking $p,m,N$ sufficiently
large therefore makes all tests pass (\Sref{Theorem}{thm:certificate-exists}). Each candidate needs finitely many
oracle queries and rational operations, so the enumeration reaches an accepted candidate in finite time. The
encodings of \Ssec{sec:alg} treat $n$ and $k$ as inputs of the same machine. The Epoch input is
admissible with disjoint source balls by Corollary~\ref{thm:main-epoch}; the Epoch machine forms the radius names
from the name of $\eps$ by total programs (the reciprocal of $1-\eps/2>\frac12$ and two square roots,
\Sref{Lemma}{lem:efd-roots-direct}) and runs the same search.
\end{proof}

\begin{remark}[The four cases]\label{rem:eff-branches}
The exact construction distinguishes the four cases of
\Sref{Construction}{cons:ham-branch}. The finite verifier has the same form in every
case: it tests rational trajectories of the proposed output. No equality test
between named capacities or decision at the threshold $c_i=C/2$ is needed.
\end{remark}

\begin{remark}[Size]\label{rem:eff-size}
A certificate with parameters $(B,e,Q,p,m,N)$ inspects $k(2^m+1)^{2n}$ node-source pairs and up to $N+1$ Euler
rows for each retained pair. The factor $4^\Theta$ in the Euler bound can require very large spatial grids and
time tables; the theorem concerns computability, not running time.
\end{remark}

\begin{remark}[The Epoch submission format]\label{rem:eff-epoch-interface}
The Epoch problem \cite{Epoch} asks for submissions as a Python function returning a PyTorch implementation of the
Hamiltonian, which a numerical verifier integrates from sampled boundary points of the source balls within fixed time
limits. We read its variables as $(Q,P)=(x,y)$, with $X_H=(\partial_PH,-\partial_QH)$, and its target as the
time-one map. The results here answer the mathematical part of the problem: Theorem~\ref{thm:intro-4d} gives the
Hamiltonian, Proposition~\ref{thm:intro-stability} the vector-field tolerance that a numerical implementation must
meet, and Theorem~\ref{thm:eff-algorithm} a uniform algorithm with finite output. A program within the verifier's
limits of ten minutes per solution and ten seconds per evaluation is not part of this paper.
\end{remark}
